\section{Erd\H{o}s Problem \#776: independent continuation}
\noindent Date: 23 September 2026. Source versions read in full: \texttt{776.tex} in \texttt{gpt\_pro\_5.4}.

\subsection*{1) OBJECTIVE AND VERIFIED BACKGROUND}
The problem asks for the threshold beyond which an antichain can occupy $n-3$ levels with $r$ sets on each level. He--Tang's revised primary paper confirms $n_0(2)=3$, $n_0(3)=8$, and $n_0(r)=2r+o(r)$. These are existing results. Below is a self-contained structural explanation for the leading linear obstruction, weaker than their sharp lower bound.
\subsection*{2) EXTREME LEVELS FORCE STARS}
Suppose $r\ge5$, $n\ge r+1$, and an antichain occupies $n-3$ levels with exactly $r$ sets on each. If it includes $r$ singletons, every larger member avoids their $r$ elements, so all occupied ranks are at most $n-r$. It then occupies at most $n-r<n-3$ ranks, a contradiction. Complementation excludes rank $n-1$ as well. Empty and full sets cannot occur with multiplicity $r>1$. Hence the occupied ranks are exactly $2,\ldots,n-2$.
Let $E$ be the family of its $r$ two-element sets and $F$ the complements of its $r$ sets of size $n-2$. Incomparability says every $e\in E$ meets every $f\in F$. If $E$ contains two disjoint edges, there are only four two-element sets meeting both, contrary to $|F|\ge5$. Thus $E$ is pairwise intersecting. A pairwise-intersecting family of at least four edges is a star: two edges $\{c,u\},\{c,v\}$ force an edge missing $c$ to be $\{u,v\}$, after which no fourth edge is possible. Similarly $F$ is a star.
Their centres must coincide. If their centres were $c\ne d$, take an edge of the first star not containing $d$; only two edges of the star at $d$ can meet it, contradicting $r\ge5$.
\subsection*{3) A SELF-CONTAINED NECESSARY SIZE BOUND}
Write $E=\{\{c,u\}:u\in U\}$ and $F=\{\{c,v\}:v\in V\}$, with $|U|=|V|=r$. Any intermediate member containing $c$ avoids all of $U$; any member missing $c$ contains all of $V$. In particular every member of rank $r-1$ contains $c$ and avoids $U$, so there are at most
\[
\binom{n-r-1}{r-2}
\]
such members. For $r+1\le n\le2r$, this is at most $r-1$, whereas rank $r-1$ must contain $r$ members. Hence necessarily $n\ge2r+1$.
This is weaker than the established threshold lower bound, but identifies the cross-intersecting-star structure responsible for the coefficient two.
\subsection*{7) FINAL}
\textbf{UNRESOLVED.} No exact threshold for general $r$ is obtained. The source's asymptotic conclusion remains an attributed result, and the argument here is a separate elementary partial analysis.
\noindent Primary source checked: He and Tang, version 2 (21 March 2026), \url{https://arxiv.org/abs/2602.09803v2}.

\subsection*{8) COMPLETION ESTIMATE}
\textbf{COMPLETION: 15\%}
This is a subjective estimate of progress toward the entire stated problem, not a probability that the conjecture or an unverified proof is correct.
